\documentclass[11pt,a4paper]{article}

\usepackage[T1]{fontenc}
\usepackage[margin=25mm]{geometry}
\usepackage{booktabs}
\usepackage{enumitem}
\usepackage[hidelinks]{hyperref}
\usepackage{xcolor}

\newcommand{\code}[1]{\texttt{\detokenize{#1}}}

\title{LIGO GNSS Processing: Implementation Modifications}
\author{Engineering Change Summary}
\date{2 August 2026}

\begin{document}
\maketitle

\section{Scope}

This document summarizes the modifications made to the LIGO GNSS processing
pipeline.  The work had two objectives: improve the structure and
maintainability of the original monolithic GNSS processing implementation,
add a base-station/rover carrier-phase processing path that produces a
floating-point RTK solution through double-difference factors, and resolve
validated carrier ambiguities to integers for a fixed RTK solution.

\section{GNSS processing redesign}

The large processing flow in \code{src/GNSS_Processing_fg.cpp} was separated
into focused operations with explicit responsibilities.  The extracted logic
includes:

\begin{itemize}[leftmargin=*]
  \item GNSS buffer clearing and initialization-window shifting;
  \item receiver clock prediction;
  \item insertion of frame state estimates into GTSAM;
  \item extraction of optimized states from the iSAM2 estimate;
  \item carrier-phase history pruning by frame number and observation time;
  \item LiDAR information-matrix normalization; and
  \item selection of the first valid constellation clock bias during GNSS
        initialization.
\end{itemize}

Repeated buffer-management and state-copying blocks were replaced by these
helpers.  Factor insertion order and factor identifiers were retained so that
the existing fixed-lag graph-deletion behavior remains compatible.  LiDAR
information normalization now detects a non-positive diagonal before scaling,
preventing division by zero and marking the LiDAR update invalid for that
epoch.

\section{Base-station observation input}

The new \code{BaseStationData} component reads surveyed base-station
coordinates and carrier observations from RINEX 3.x observation data, with
RINEX 3.02 verified explicitly.  Its input may be a plain observation file, a
ZIP archive containing an observation member, or a directory of files and
archives.  ZIP members are read in place through \code{libzip}; no manual
extraction or CSV conversion is required.  The base antenna position is read
in Earth-Centered, Earth-Fixed (ECEF) coordinates from:

\begin{verbatim}
APPROX POSITION XYZ
\end{verbatim}

The reader assembles multiline \code{SYS / # / OBS TYPES} header records,
reads primary/secondary/L5 carrier phase and LLI fields, applies GLONASS frequency channels,
deduplicates multiple tracking codes on one satellite/frequency, and handles
special RINEX event records.  GPS and Galileo bands 1 and 5, BeiDou B1I
(RINEX 3.02 band 1, with band 2 also accepted) and B2a, and GLONASS L1 are
mapped to their physical carrier frequencies.  UTC and BDT observation times
are converted to GPST before epochs are stored and matched to the rover time
within a configurable tolerance.

RINEX LLI and SSI fields do not encode carrier standard deviation, so the
reader assigns the configured \code{base_carrier_std_cycles} value.  It
rejects malformed headers, absent ECEF coordinates, inconsistent positions
across files, and inputs with no supported carrier observations.  The active
configuration points to \code{Data/2026-Aug-02_001813}, which contains the
08:00 and 09:00 GPST zipped HKSC observation files for 21 March 2025.  These
epochs overlap the rover bag in \code{Dataset}; the enclosing directory date
is only the local download name.

\section{Double-difference carrier factors}

For every synchronized rover/base epoch, the processing pipeline performs the
following operations:

\begin{enumerate}[leftmargin=*]
  \item Match rover and base observations by satellite and carrier frequency.
  \item Retain the primary band and, when enabled, GPS L5, Galileo E5a, and
        BeiDou B2a observations at 1176.45 MHz.
  \item Group observations by GNSS constellation and signal band.
  \item Select the lowest-noise carrier observation in each
        constellation/band group as its reference satellite.
  \item Form the carrier measurement double difference
        \[
          \Delta\nabla L =
          (L_{r}^{s}-L_{b}^{s})-(L_{r}^{q}-L_{b}^{q}),
        \]
        where $r$ and $b$ denote rover and base receivers, $s$ is the subject
        satellite, and $q$ is the same-constellation reference satellite.
  \item Add a geometric range double-difference factor and a persistent,
        real-valued ambiguity state to the iSAM2 graph.
\end{enumerate}

Two factor variants were implemented in
\code{include/gnss_factor/gnss_double_diff_carrier_factor.hpp}.  The first
uses the local LiDAR/IMU position together with the estimated local-to-ECEF
alignment.  The second directly uses the ECEF navigation state when LiDAR is
disabled.  Both factors include analytical Jacobians for the connected pose,
position, alignment, and ambiguity variables.

The ambiguity associated with a satellite/reference/band tuple is initialized from
the difference between measured and predicted geometry.  It remains a
continuous GTSAM \code{Vector1} state across epochs.  A new ambiguity is
created after a receiver or base loss-of-lock indication, an observation gap,
or a change of reference pair. L1 and L5 therefore have independent ambiguity
arcs and wavelengths. Until integer validation succeeds, these
states provide the floating-point RTK solution.

\section{LAMBDA integer ambiguity resolution}

The new \code{LambdaAmbiguityResolver} implements integer least-squares
decorrelation and search.  Ambiguities and their full joint marginal
covariance are converted from metres to cycles before resolution.  Only
same-frequency double differences are eligible; this excludes combinations
whose metre-valued ambiguity cannot be represented by one integer wavelength.
All eligible primary-band and L5 ambiguities enter the same joint covariance
and integer search; no ionosphere-free carrier combination is used because
that would destroy the direct integer nature of the original ambiguities.

Integer fixing is guarded by a minimum number of ambiguities, a continuous
lock-epoch requirement, a maximum float standard deviation, and the ratio of
the second-best to best LAMBDA squared residual.  Partial ambiguity resolution
removes the highest-variance ambiguity and retries when the complete set is
not sufficiently precise or does not pass validation.  Failure at any gate
leaves the float graph unchanged.

After a successful ratio test, each accepted integer is converted back to
metres and inserted as a tight ambiguity prior.  iSAM2 is updated a second
time in the same epoch, so the constrained navigation state is returned
immediately.  The constraint is held for the continuous ambiguity arc; a
receiver/base loss-of-lock flag or observation gap starts a new float
ambiguity instead of reusing the old integer.

\section{Configuration}

The following parameters were added to the GNSS section of
\code{config/avia.yaml}:

\begin{center}
\begin{tabular}{@{}lll@{}}
\toprule
Parameter & Default & Purpose \\
\midrule
\code{use_double_differences} & true & Enable base/rover carrier factors \\
\code{use_secondary} & true & Add GPS L2, Galileo E5b, and BeiDou B2 \\
\code{use_l5} & true & Add supported 1176.45 MHz carrier factors \\
\code{base_station_file} & HKSC directory & RINEX file, ZIP, or directory under \code{Data} \\
\code{base_epoch_tolerance} & 0.05 s & Maximum rover/base epoch offset \\
\code{rtk_ambiguity_gap_tolerance} & 1.5 s & Continuous-arc gate for 1 Hz base data \\
\code{base_carrier_std_cycles} & 0.01 cycles & Assumed RINEX phase uncertainty \\
\code{double_difference_sigma} & 0.02 m & Carrier factor standard deviation \\
\code{ambiguity_prior_sigma} & 100 m & Loose float-ambiguity prior \\
\code{enable_integer_fixing} & true & Enable validated LAMBDA fixing \\
\code{lambda_min_ambiguities} & 3 & Smallest partial-fix subset for this dataset \\
\code{lambda_min_lock_epochs} & 10 & Required continuous ambiguity arc \\
\code{lambda_ratio_threshold} & 3.0 & Integer candidate ratio gate \\
\code{lambda_max_std_cycles} & 0.25 cycles & Float precision gate \\
\code{fixed_ambiguity_sigma_cycles} & 0.001 cycles & Fix-and-hold prior sigma \\
\code{rtk_debug} & true & Emit structured RTK diagnostics \\
\code{rtk_debug_epoch_interval} & 10 epochs & Period for repeated float/DD reports \\
\bottomrule
\end{tabular}
\end{center}

Double differences are enabled for the synchronized rover bag and the
08:00--10:00 GPST HKSC base observations.

\section{RTK runtime diagnostics (2 August 2026)}

Structured log messages now make the RTK state explicit.  \code{[RTK-DD]}
reports synchronized epochs, reference satellites per constellation/band,
matched candidates, factors added, L5 factors, new ambiguity arcs, cycle-slip
resets, missing base epochs, zero-factor epochs, and cumulative DD totals.  A new arc also prints its measured DD,
geometric DD, and initial float ambiguity in metres and cycles.

\code{[RTK-FLOAT]} prints active float/fixed counts, lock progress, each
eligible ambiguity estimate in cycles, and its marginal standard deviation.
\code{[RTK-LAMBDA]} prints the attempt number, precision rejections, solver
rejections, the best and second-best squared norms, ratio-test decisions, and
each accepted integer with its float-to-fixed correction.  A successful fix is
unambiguous: the log contains \code{FIX SUCCESS} and the public status is
\code{FIXED}.  Statuses such as \code{FLOAT_WAITING_FOR_LOCK},
\code{FLOAT_REJECTED_PRECISION}, and \code{FLOAT_REJECTED_RATIO} identify why
the solution remains float.  Public counters expose attempts, successful
fixes, and current float/fixed ambiguity counts to debuggers.

\section{RTK continuity repair and replay (2 August 2026)}

The ambiguity continuity gate is now independent of the 10~Hz rover sampling
period.  A 1.5~s RTK gap tolerance preserves arcs across consecutive 1~Hz
RINEX epochs.  RINEX LLI is interpreted as a bit mask: bits 0 and 1 invalidate
an integer arc, while Galileo bit 2 only marks BOC tracking and no longer
causes a false cycle-slip reset.  Reference satellites are held while still
observed, and LAMBDA is invoked only after an epoch actually adds DD factors.

The reader and DD classifier now include the second frequencies present in the
bag: GPS L2, Galileo E5b, and BeiDou B2.  Real-data compatibility consequently
increased from 2975 to 4449 matched carriers and from 2100 to 3011 potential
DD factors; the bag contains 31825 secondary-frequency observations and no L5.

A real ROS replay confirmed that ambiguity keys are reused across consecutive
base epochs (for example the second synchronized epoch reported zero new or
reset arcs).  At lock epoch 10, LAMBDA executed with three eligible float
ambiguities.  Their marginal standard deviations were approximately
0.85--1.41 cycles, above the configured 0.25-cycle gate, so the integer fix was
correctly rejected as \code{FLOAT_REJECTED_PRECISION}.  The processing defect
is repaired, but this replay does not claim a validated fixed solution.

\section{Build and verification}

The GNSS processor header and implementation now contain structured source
annotations for the class lifecycle, graph-key convention, initialization
branches, fixed-lag update, state import/export, satellite preparation, six
factor-assembly stages, multi-frequency ambiguity arcs, and LAMBDA validation.  The
comments describe ownership, frames, and failure behavior rather than merely
restating individual C++ expressions.  A separate pipeline reference is
provided in \code{paper/GNSS_Processing_fg_pipeline.tex}.
The LAMBDA resolver is likewise annotated with its matrix convention,
unimodular transform, integer Gauss operations, conditional-variance
permutations, depth-first zig-zag enumeration, pruning radius, candidate
ordering, back-transformation, and float-fallback failure behavior.

The build definition now compiles \code{src/BaseStationData.cpp} and
\code{src/LambdaAmbiguityResolver.cpp} into \code{ligo_mapping}.  The
base reader and its test executable link \code{libzip}.  The
package-level CMake configuration no longer
overwrites a caller-supplied build type, which also resolves a mismatch between
the package and catkin's bundled GoogleTest library.

The complete ROS executable was built and linked successfully with a
single-job Release build.  Base-file parsing remains in
\code{test/test_base_station_data.cpp}.  A separate RTK core suite in
\code{test/test_rtk_core.cpp} verifies both double-difference factor variants,
including zero residuals and every analytical Jacobian against central finite
differences.  It also verifies a known correlated LAMBDA solution, candidate
ordering, integer-translation invariance, invalid-input rejection, and 40
deterministic covariance cases against exhaustive integer search.

The RTK suite also verifies primary/secondary/L5 classification, LLI bit-mask
handling, RTK gap continuity, frequency-specific base
matching, disabling L5, and rejection of unsupported GLONASS L5.  The base
reader suite verifies a compact multiline RINEX 3.02 fixture and directly
loads the two actual HKSC ZIP archives, including multi-frequency extraction, surveyed
ECEF coordinates, aggregation across both hours, epoch matching, and a flag-4
receiver event.  A separate diagnostic test prints per-constellation primary
and L5/E5a satellite counts from the first epoch of the actual HKSC ZIP.  A
rosbag compatibility test scans all 3040 rover epochs and reports 304
synchronized epochs, 4449 matched carrier observations, and 3011 potential DD
factors.  The four LIGO test executables reported fifteen passing tests.  The complete catkin result
collector, including the workspace's generated test records, reported:

\begin{center}
  \textbf{30 tests, 0 errors, 0 failures, 0 skipped.}
\end{center}

The bag contains no 1176.45 MHz rover observations, so the real-data path can
exercise L1 but not L5.  A headless replay loads and initializes the IMU, but
does not reach GNSS graph processing because the bag publishes
\code{livox_ros_driver2/CustomMsg} while this repository consumes
\code{livox_ros_driver/CustomMsg}; the experimental no-LiDAR loop also does
not dispatch the GNSS callback.  A message adapter or driver-2 port is needed
for a complete factor-graph replay and runtime accuracy evaluation.

\end{document}
